% ============================================================
% DEMO 2 — Mathematical notation for statistics
% Live demo: show custom commands, environments, aligned eq
%
% Compile with:  xelatex math.tex
% ============================================================

\documentclass[12pt, a4paper]{article}

\usepackage{amsmath, amssymb, amsthm}
\usepackage{hyperref}

% ----- Custom commands (define once, use everywhere) -----
\newcommand{\E}{\mathbb{E}}
\newcommand{\Var}{\mathrm{Var}}
\newcommand{\Cov}{\mathrm{Cov}}
\newcommand{\iid}{\overset{\text{iid}}{\sim}}
\newcommand{\norm}[1]{\left\lVert #1 \right\rVert}
\newcommand{\abs}[1]{\left\lvert #1 \right\rvert}
\newcommand{\given}{\mid}

% ----- Theorem environments -----
\newtheorem{theorem}{Theorem}[section]
\newtheorem{lemma}[theorem]{Lemma}
\newtheorem{corollary}[theorem]{Corollary}
\theoremstyle{definition}
\newtheorem{definition}{Definition}[section]
\theoremstyle{remark}
\newtheorem{remark}{Remark}

\title{Mathematical Notation for Statistics}
\author{Imad El Badisy}
\date{\today}

% ============================================================
\begin{document}

\maketitle
\tableofcontents
\newpage

% ----- Scalars, vectors, matrices -----
\section{Notation Conventions}

We use the following conventions throughout:

\begin{itemize}
  \item Scalars: $x$, $\mu$, $\sigma^2$ (lowercase italic)
  \item Vectors: $\mathbf{x}$, $\boldsymbol{\beta}$ (bold lowercase)
  \item Matrices: $\mathbf{X}$, $\boldsymbol{\Sigma}$ (bold uppercase)
  \item Sets: $\mathcal{X}$, $\mathcal{F}$ (calligraphic)
  \item Random variables: $X$, $Y$ (uppercase italic)
\end{itemize}

% ----- Probability -----
\section{Probability and Expectation}

Let $X$ and $Y$ be random variables on a probability
space $(\Omega, \mathcal{F}, \Pr)$.

\begin{align}
  \E[X]        &= \int x \, f(x) \, dx                    \\
  \Var(X)      &= \E\left[(X - \E[X])^2\right]            \\
  \Cov(X, Y)   &= \E\left[(X - \E[X])(Y - \E[Y])\right]
\end{align}

We write $X_1, \ldots, X_n \iid N(\mu, \sigma^2)$ to mean
that the $X_i$ are independent and identically distributed.

% ----- Statistical models -----
\section{Statistical Models}

\subsection{Linear Model}

\begin{equation}
  \mathbf{y} = \mathbf{X}\boldsymbol{\beta} + \boldsymbol{\varepsilon},
  \qquad
  \boldsymbol{\varepsilon} \sim N(\mathbf{0}, \sigma^2 \mathbf{I}_n)
  \label{eq:lm}
\end{equation}

The OLS estimator minimizes
$\norm{\mathbf{y} - \mathbf{X}\boldsymbol{\beta}}^2$:

\begin{equation}
  \hat{\boldsymbol{\beta}}
    = (\mathbf{X}^\top \mathbf{X})^{-1} \mathbf{X}^\top \mathbf{y}
\end{equation}

\subsection{Logistic Regression}

For a binary outcome $Y_i \in \{0,1\}$:

\begin{equation}
  \log \frac{p_i}{1 - p_i}
    = \mathbf{x}_i^\top \boldsymbol{\beta},
  \qquad
  p_i = \Pr(Y_i = 1 \given \mathbf{x}_i)
\end{equation}

\subsection{Cox Proportional Hazards}

\begin{equation}
  h(t \given \mathbf{x})
    = h_0(t) \exp\!\left(\mathbf{x}^\top \boldsymbol{\beta}\right)
  \label{eq:cox}
\end{equation}

where $h_0(t)$ is the unspecified baseline hazard.

% ----- Likelihood -----
\section{Likelihood and Information}

The log-likelihood is:

\begin{equation}
  \ell(\boldsymbol{\theta})
    = \sum_{i=1}^{n} \log f(y_i \given \mathbf{x}_i, \boldsymbol{\theta})
\end{equation}

The Fisher information matrix is:

\begin{equation}
  \mathcal{I}(\boldsymbol{\theta})
    = -\E\!\left[
        \frac{\partial^2 \ell(\boldsymbol{\theta})}
             {\partial \boldsymbol{\theta} \, \partial \boldsymbol{\theta}^\top}
      \right]
\end{equation}

% ----- Theorems -----
\section{Theorem Environments}

\begin{definition}[Unbiased estimator]
  An estimator $\hat{\theta}$ is \emph{unbiased} if
  $\E[\hat{\theta}] = \theta$ for all $\theta$.
\end{definition}

\begin{theorem}[Cramér-Rao lower bound]
  \label{thm:crlb}
  For any unbiased estimator $\hat{\theta}$ of a scalar
  parameter $\theta$,
  \begin{equation}
    \Var(\hat{\theta}) \geq \frac{1}{\mathcal{I}(\theta)}
  \end{equation}
  where $\mathcal{I}(\theta)$ is the Fisher information.
\end{theorem}

\begin{proof}
  Apply the Cauchy-Schwarz inequality to
  $\Cov\!\left(\hat{\theta}, \partial_\theta \log f(X;\theta)\right)$.
  The score has mean zero and variance $\mathcal{I}(\theta)$,
  which gives the result.
\end{proof}

\begin{corollary}
  The OLS estimator in \eqref{eq:lm} is the minimum-variance
  unbiased linear estimator (Gauss-Markov theorem).
\end{corollary}

\end{document}
